\documentclass[10pt,letterpaper]{article}
\usepackage[margin=0.7in,headheight=14pt]{geometry}
\usepackage[T1]{fontenc}
\usepackage{lmodern,amsmath,amssymb,xcolor,fancyhdr,microtype}
\usepackage{titlesec}
\usepackage[hidelinks]{hyperref}
\definecolor{navy}{HTML}{192B4D}
\pagestyle{fancy}\fancyhf{}
\fancyhead[L]{\small MF 731 | 07: Risk Measures}
\fancyhead[R]{\small Homework Solutions}
\fancyfoot[C]{\small\thepage}
\setlength{\parindent}{0pt}
\setlength{\parskip}{5pt}
\titlespacing*{\section}{0pt}{12pt}{5pt}
\hypersetup{pdftitle={MF 731 - Entropic Risk Measure: Solutions},
 pdfsubject={Cash-additivity, monotonicity, convexity and failure of positive homogeneity}}
\begin{document}
\setlength{\abovedisplayskip}{7pt}
\setlength{\belowdisplayskip}{7pt}
\setlength{\abovedisplayshortskip}{4pt}
\setlength{\belowdisplayshortskip}{5pt}
{\LARGE\bfseries\textcolor{navy}{Entropic Risk Measure}}

For a loss $L$ and fixed $\gamma>0$, define
\[
\rho(L)=\frac{1}{\gamma}\ln\mathbb E[e^{\gamma L}].
\]
We consider losses for which this expectation is finite.

\section*{1. Cash-additivity and monotonicity}
For any constant $c\in\mathbb R$,
\[
\rho(L+c)=\frac{1}{\gamma}\ln\left(e^{\gamma c}\mathbb E[e^{\gamma L}]\right)
=c+\frac{1}{\gamma}\ln\mathbb E[e^{\gamma L}]
=\boxed{\rho(L)+c.}
\]
Thus adding a certain loss $c$ increases the required capital by $c$.
Equivalently, adding cash $c$ reduces the loss to $L-c$ and gives
$\rho(L-c)=\rho(L)-c$.

If $L_1\le L_2$ almost surely, then $e^{\gamma L_1}\le e^{\gamma L_2}$ almost surely.
Taking expectations and using that $\ln$ is increasing and $\gamma>0$ gives
\[
\boxed{\rho(L_1)\le\rho(L_2).}
\]

\section*{2. Convexity using H\"older's inequality}
Let $0<\theta<1$, and set $p=1/\theta$ and $q=1/(1-\theta)$, so
$1/p+1/q=1$. H\"older's inequality gives
\begin{align*}
\mathbb E\!\left[e^{\gamma(\theta L_1+(1-\theta)L_2)}\right]
&=\mathbb E\!\left[(e^{\gamma L_1})^\theta
                         (e^{\gamma L_2})^{1-\theta}\right]\\
&\le\left(\mathbb E[e^{\gamma L_1}]\right)^\theta
      \left(\mathbb E[e^{\gamma L_2}]\right)^{1-\theta}.
\end{align*}
Taking logarithms and dividing by $\gamma$ yields
\[
\boxed{\rho(\theta L_1+(1-\theta)L_2)
\le\theta\rho(L_1)+(1-\theta)\rho(L_2).}
\]
The cases $\theta=0$ and $\theta=1$ are equalities. Together with Question 1,
this proves that $\rho$ is a convex risk measure.

\section*{3. Failure of coherence}
Let $L\sim N(0,1)$ and $\lambda>0$. Completing the square gives
\begin{align*}
\mathbb E[e^{\gamma\lambda L}]
&=\frac{1}{\sqrt{2\pi}}\int_{-\infty}^{\infty}
       e^{\gamma\lambda x-x^2/2}\,dx\\
&=e^{\gamma^2\lambda^2/2}\frac{1}{\sqrt{2\pi}}
  \int_{-\infty}^{\infty}e^{-(x-\gamma\lambda)^2/2}\,dx
=e^{\gamma^2\lambda^2/2}.
\end{align*}
Therefore
\[
\boxed{\rho(\lambda L)=\frac{\gamma\lambda^2}{2},}
\qquad
\lambda\rho(L)=\frac{\gamma\lambda}{2}.
\]
For example, $\rho(2L)=2\gamma$ whereas $2\rho(L)=\gamma$.
Positive homogeneity fails, so the entropic risk measure is \textbf{not coherent}.
\end{document}
